% ECONOMICS_PROBLEM: {"id": "O31-155", "title": "Financing innovation", "jel_code": "O31", "source_id": "OP-0267"}
% !TeX program = xelatex
\documentclass[11pt,a4paper]{article}
\usepackage[margin=23mm]{geometry}
\usepackage{fontspec}
\setmainfont{Latin Modern Roman}
\usepackage{amsmath,amssymb,mathtools}
\usepackage{xcolor,enumitem,booktabs,longtable,array,xurl,hyperref}
\definecolor{muted}{HTML}{53636C}
\hypersetup{colorlinks=true,urlcolor=blue,linkcolor=blue}
\setlength{\parindent}{0pt}
\setlength{\parskip}{5pt}
\newcommand{\R}{\mathbb R}
\newcommand{\E}{\mathbb E}
\newcommand{\N}{\mathbb N}
\newcommand{\ind}{\mathbf 1}
\newcommand{\Var}{\operatorname{Var}}
\newcommand{\Cov}{\operatorname{Cov}}
\newcommand{\supp}{\operatorname{supp}}
\DeclareMathOperator*{\argmin}{arg\,min}
\DeclareMathOperator*{\argmax}{arg\,max}
\newcommand{\entrymeta}[3]{{\sffamily\small\color{muted}\textbf{\texttt{#1}}\quad|\quad #2\par #3\par}}
\newcommand{\Problemref}[1]{\texttt{#1}}
\newcommand{\JELref}[2]{\texttt{#2} (JEL \texttt{#1})}
\begin{document}
\section*{Financing innovation}
\textbf{Problem ID:} \texttt{O31-155}\quad \textbf{JEL:} \texttt{O31}\par
% BEGIN ECONOMICS STATEMENT
\label{problem:OP-0267}
{\sffamily\small\color{muted}Original subject group: Growth and productivity\par}
\label{definitions:problem:30007548}
\textbf{Audit classification:} Explicit published open question. Finance and innovation are expressly open research directions; hidden-quality public finance mechanism is editorial and strategic. Verification concerns the cited version, not first priority or proof that this exact editorial specification remains unsolved on October 8, 2026.
\subsection*{Definitions and precise research target}
\paragraph{Editorial mathematical specification.} Consider young firms with privately known research quality \(q\), wealth \(w\), and observable signals \(s\). A project requires funding \(k\), produces an uncertain innovation with specified probability \(p(q,k,e)\), and depends on unobserved effort \(e\). Private cash flows differ from social innovation value because of knowledge spillovers and business stealing. Investors choose contracts contingent on observable outcomes; public policy can provide a bounded guarantee or co-investment budget. Let \(\mathcal C\) be a stated feasible contract family.
Find financing and public-support rules maximizing expected social innovation surplus net of research resources, screening costs, and the financing cost of public funds. Require lender participation, entrepreneur participation, truthful quality revelation where reports are used, and effort incentive compatibility. For each type, a financed project must satisfy the declared equilibrium contract constraints rather than an assumed social-value threshold.
Determine when guarantees or equity co-investment expand valuable radical innovation and when they mainly fund weak projects or displace private finance. Identify the quality distribution, effort response, and spillover parameters using funding discontinuities, contract variation, and independently measured innovation outcomes. Compare policies under common fiscal budgets and report robustness to private-information misspecification. The target is a scoped financing mechanism for young-firm innovation; it does not claim that financial constraints or entrepreneurship have no established theoretical treatment.

For a one-project benchmark take compact quality and effort sets, \(0\le p(q,k,e)\le1\), success cash flow \(X(q)\ge0\), social success value \(V(q)\), and effort cost \(c(e,q)\). A contract specifies financing \(k\), entrepreneur repayment \(R_0,R_1\) by failure/success, and public guarantees \(G_0,G_1\), all contingent only on verifiable outcomes and reports. Entrepreneur expected utility is \(p(X-R_1)-(1-p)R_0-c(e,q)\); specify limited liability, outside utility, and wealth contribution. Effort must maximize this expression for the assigned contract, and truthful quality reporting must dominate every alternative report with effort reoptimized. Lender participation includes \(p(R_1+G_1)+(1-p)(R_0+G_0)\) relative to actual funded resources. Social transfers are not real costs except through their specified financing distortion. The optimal financing mechanism involves strategic private information and hence overlaps game theory.
\subsection*{Variants and assumption changes}
The following are \emph{editorial variants}, not additional published open conjectures.
\begin{enumerate}
\item \emph{Hidden effort only.} Let quality be public and vary guarantee coverage within a fixed contract family. Characterize the tradeoff between relaxing lender risk and weakening entrepreneurial effort incentives under limited liability.
\item \emph{Hidden quality and screening.} Permit noisy publicly purchased screening and an incentive-compatible menu of equity co-investment contracts. Compare the extra valuable projects with displacement of private funding under the same public budget.
\end{enumerate}
\subsection*{References and source audit}
\begin{itemize}
\item Ufuk Akcigit. \emph{Creative destruction in economic growth}. 2026. Locator: Section7, Open questions, item1. Primary text checked. \url{https://onlinelibrary.wiley.com/doi/10.1111/sjoe.70037}.
\end{itemize}
Finance and innovation are expressly open research directions; hidden-quality public finance mechanism is editorial and strategic. This formulation contains strategic incentives and therefore overlaps game theory.
\begin{enumerate}\raggedright
\item\hypertarget{definitions:source:30007548:1}{} Ufuk Akcigit. 2026. \emph{Creative destruction in economic growth}. Section7, Open questions, item1.
\url{https://onlinelibrary.wiley.com/doi/10.1111/sjoe.70037}.
DOI: \url{https://doi.org/10.1111/sjoe.70037}.
\end{enumerate}

\subsection*{Common conventions: Source mathematical conventions}

These conventions apply to each entry unless explicitly replaced.

\textbf{Models and identification.} A primitive model \(m\) specifies all probability laws, feasible choices, information, equilibrium and measurement rules used by its target. The admissible class \(\mathcal M\) is fixed independently of outcomes. If \(P_O(m)\) is its observable probability law and \(\tau(m)\) the target, its identified set at observable law \(P\) is
\[
 \mathcal I_\tau(P)=\{\tau(m):m\in\mathcal M,\ P_O(m)=P\}.
\]
Point identification means that this set is a singleton. An empty set signals incompatibility of data and assumptions. Coordinate bounds need not describe the joint set. Bounds are sharp only if every retained target is attainable or arbitrarily approachable by compatible models. Finite bounds require explicit restrictions on unobserved quantities; unrestricted confounding, hidden wealth, extrapolation or arbitrary equilibrium selection can yield uninformative sets.

\textbf{Causal interventions.} A potential outcome \(Y_i(p)\) is the outcome under a complete policy regime \(p\), including timing, financing, information and market spillovers. The target population and units are fixed before treatment. With interference, use the complete assignment vector or a justified exposure mapping; individual no-interference notation alone cannot identify an economy-wide intervention. A difference of mean potential outcomes does not require the cross-world joint distribution. Conditioning on post-treatment migration, survival, enrollment or employment changes the estimand and needs separate justification.

\textbf{Statistical statements.} An estimator, confidence set, forecast or test specifies a sampling law, model class and loss/coverage criterion. Uniform \((1-\alpha)\)-coverage means \(\inf_{m\in\mathcal M}\Pr_m\{\tau(m)\in C_n\}\geq1-\alpha\), or an explicitly asymptotic version. The whole parameter space trivially has coverage, so informativeness requires a stated width, risk or power objective. A forecasting improvement is relative to a named benchmark, information set and evaluation population.

\textbf{Equilibrium and optimization.} An equilibrium comprises feasible plans, optimality, beliefs where needed, budget constraints and all market/resource clearing. The primitives or a stated benchmark define each correspondence \(\mathcal E(p;m)\). Nonemptiness is assumed or proved, never inferred just from compact policy sets. A selected-equilibrium target uses a declared selection rule; a robust target quantifies over all permitted equilibria. Use suprema or infima unless attainment is established. An \(\varepsilon\)-optimal policy is feasible and within \(\varepsilon\) of the finite optimum value. Report an empty feasible set separately.

\textbf{Welfare and accounting.} Normative utility, interpersonal normalization, social weights, numeraire, discounting and the use of public receipts are primitives. Behavioral utility need not coincide with the welfare criterion. Household welfare includes ownership income and rents once; adding profits again to compensating variation generally double-counts them. A monetary equivalent is defined by an explicit transfer and price convention, with monotonicity and range conditions. Average effects, mechanism contributions, transfers and welfare losses are different targets. Infinite sums require convergence and dynamic feasible plans require terminal or no-Ponzi conditions.

\textbf{Source versions.} A working-paper date, revision date and journal publication year are distinguished. A DOI for a journal version is not automatically the DOI of an earlier manuscript. Locators refer to the stated version and printed pages where specified. A metadata-only check does not verify a claim about a full-text conclusion. Every entry's source subsection narrows the provenance claim accordingly.
% END ECONOMICS STATEMENT
\end{document}
